\section{Finite presentations of measure-valued filters}\label{sec:v7-measure}
The preceding sharp example has a separated tree geometry. We next give
a constructive approximation theorem requiring no such separation and
no finite-dimensional exact posterior coordinate.

\begin{lemma}[A finite Wasserstein presentation]\label{lem:v7-wasserstein-net}
For $n\ge1$ let $\mathcal C_n$ consist of empirical probability measures
with $n$ equally weighted atoms on $\{0,1/n,\ldots,1\}$. Then
\begin{equation}
 |\mathcal C_n|=\binom{2n}{n}\le4^n,
 \qquad\sup_{\pi\in\mathcal P([0,1])}\Wone(\pi,Q_n\pi)
                              \le\frac{3}{2n}         \label{eq:v7-measure-net}
\end{equation}
for a specified Borel quantizer $Q_n$.
\end{lemma}
\begin{proof}
Take the generalized quantiles of $\pi$ at $(j-1/2)/n$, $1\le j\le n$,
and round them to nearest grid points, breaking ties toward zero.
Quantiles are Borel functions of a probability law. If $q$ is its
monotone quantile function, the integral of the error from replacing
$q$ by its midpoint value on each of the $n$ equal subintervals is at
most $1/n$: bound each integral by its interval length times the
oscillation and sum the monotone variations, at most one. Rounding
adds at most $1/(2n)$. In one dimension $\Wone$ is the $L^1$ distance
of quantile functions, giving the bound. The possible sorted lists are
multisets of size $n$ from $n+1$ points, counted by $\binom{2n}{n}$.
The binomial theorem bounds this by $4^n$.
\end{proof}

\begin{theorem}[A charged measure-valued recursive approximation]
\label{thm:v7-filter}
Let the actual filter obey $\pi_t=\Phi_{y_t,a_t}(\pi_{t-1})$ on
$\mathcal P([0,1])$, and suppose every allowed update is
$\rho$-contractive in $\Wone$, with $\rho<1$. Suppose an implementable
update has $\Wone$ defect at most $\eta$. The initialized register
\[
 \widehat\pi_t=Q_n\widetilde\Phi_{y_t,a_t}(\widehat\pi_{t-1}),
              \qquad\widehat\pi_0\in\mathcal C_n
\]
has at most $\binom{2n}{n}$ persistent values and satisfies, for every
actual observation path,
\begin{equation}
 \Wone(\pi_t,\widehat\pi_t)
 \le\rho^t+\frac{3/(2n)+\eta}{1-\rho}.                \label{eq:v7-measure-filter}
\end{equation}
This bounds the error of every $1$-Lipschitz test expectation, uniformly
over the choice of test. Thus for $M\ge4$, $n=\lfloor\log M/\log4\rfloor$
gives error $O((\log M)^{-1}+\eta)$ and squared error
$O((\log M)^{-2}+\eta^2)$ after the initialization term. If these test
expectations are conditional means, the latter is a Bayes excess bound.
Report-law comparison additionally requires the report Lipschitz
hypothesis of Theorem~\ref{thm:v7-transport}; it is not automatic from
$\Wone$ contraction.
\end{theorem}
\begin{proof}
Apply Theorem~\ref{thm:v7-transport} with the true metric state
$S=\mathcal P([0,1])$, radius $3/(2n)$, and diameter one. The dual
formula for $\Wone$ gives the assertion for every test; independent
selection of a terminal test does not enlarge its bound. The codebook
count and the logarithmic substitution are exact. Only an index of a
codebook law persists between updates, not its exact unquantized
posterior or a numerical residual.
\end{proof}

\subsection{A nonlinear physical realization}
Let $B_t$ be independent fair bits and define
\begin{equation}
 h_0(x)=x/5+x^2/50,\qquad
 h_1(x)=3/5+x/5+x^2/50,\qquad X_t=h_{B_t}(X_{t-1}).      \label{eq:v7-nonlinear-ifs}
\end{equation}
The observation $Y_t$ reports $B_t$ correctly with probability $3/4$,
independently at each time. A permitted terminal probe $g:[0,1]\to[0,1]$
returns a Bernoulli variable of probability $g(X_t)$; the tested family
is $1$-Lipschitz probes and their bounded finite executable continuations.
The actual update-report instrument on measures is
\[
 \mathsf T_y\pi=\tfrac12\{\tfrac34(h_y)_\#\pi+
                         \tfrac14(h_{1-y})_\#\pi\}.
\]
Its mass is $1/2$, independent of the old state. The normalized filter is
therefore
\begin{equation}
 \Phi_y\pi=\tfrac34(h_y)_\#\pi+\tfrac14(h_{1-y})_\#\pi. 
                                                               \label{eq:v7-nonlinear-filter}
\end{equation}

\begin{corollary}[Nonlinear posterior, explicit register]
\label{cor:v7-nonlinear}
The actual measure-valued filter \eqref{eq:v7-nonlinear-filter} is
$6/25$-contractive in $\Wone$. Theorem~\ref{thm:v7-filter} applies with
$\eta=0$ to an exact rational implementation on $\mathcal C_n$.
The model has no fixed-step nonzero common minorizer. For every integer
$D\ge1$, the first $D$ polynomial moments do not determine their own
exact next-step update on the space of all prior laws.
\end{corollary}
\begin{proof}
The two images lie respectively in $[0,11/50]$ and $[3/5,41/50]$.
Each derivative is between $1/5$ and $6/25$. Couple two old laws
optimally and use the same branch draw in their mixtures; this gives
$\Wone(\Phi_y\pi,\Phi_y\pi')\le(6/25)\Wone(\pi,\pi')$.
For a codebook input, \eqref{eq:v7-nonlinear-filter} has at most $2n$
atoms, at rational points with denominator dividing $50n^2$, with
weights multiples of $1/(4n)$. Sorted cumulative weights give the exact
quantiles, which are rounded back to the fixed grid. Thus no numerical
limit is being hidden in $\eta=0$. The transient list of at most $2n$
atoms is discarded; its entries and weights have $O(\log n)$-bit
rational descriptions. The designated persistent object is the single
codebook index. Approximate arithmetic instead obeys the separately
charged $\eta$ bound of the theorem.

After $k$ hidden transitions the law from a fixed $x$ has finite support
$\{h_v(x):|v|=k\}$. For a second initial point $x'$, each equality
$h_v(x')=h_u(x)$ excludes at most one $x'$, since $h_v$ is strictly
increasing. Choose $x'$ outside this finite set. The two supports are
disjoint, so a measure dominated by all $k$-step laws must be zero.

For the moment assertion, the $D$th updated moment integrates
$\tfrac34 h_y(x)^D+\tfrac14 h_{1-y}(x)^D$, a polynomial of degree $2D$
with nonzero leading coefficient $50^{-D}$. It is not in the span of
$1,x,\ldots,x^D$. Linear algebra on finitely many distinct points
therefore supplies a signed measure annihilating this span but not the
new polynomial. Its positive and negative parts have equal mass;
normalizing them yields two prior probabilities with identical first
$D$ moments and different updated $D$th moment. The argument is about
ordinary finite moment closure, not the impossibility of a Borel
encoding of measures by real numbers.
\end{proof}

The quotient here is genuinely the posterior law for the specified
probe family: bounded Lipschitz functions separate probability measures
on a compact interval, and constants do not affect that separation.
Updates respect this quotient by \eqref{eq:v7-nonlinear-filter}. The
advancing reports have the exact state-independent fair law, while a
terminal Bernoulli probe has total-variation defect at most
$\Wone(\pi_t,\widehat\pi_t)$. Thus the same finite presentation defines
a causal experiment simulation with a controlled terminal error and an
exactly charged register. Multiple probes would update the posterior
again; no claim for an unverified enlarged probe protocol is implicit.
This example connects positive instruments, predictive quotient,
recursive approximation and resource transformation within one model.
